% ============================================================
% Proyecto Pedagógico — páginas 2 a 5 del original, transcriptas
% a LaTeX tal cual (texto literal, mismo formato y cortes de página).
% Requiere en el preámbulo:
%   \usepackage{xcolor} \usepackage{fancyhdr} \usepackage{enumitem}
%   \usepackage{microtype} \usepackage{geometry}
% Va a continuación de la carátula: % ============================================================
% Carátula — reemplaza el bloque \begin{titlepage}...\end{titlepage}
% Requiere en el preámbulo: \usepackage{graphicx} y \usepackage{tikz}
% Logos en la carpeta logos/ (logo-impa.pdf y logo-pba.pdf)
% ============================================================
\begin{titlepage}
% Logo PBA en el margen superior derecho (posición absoluta en la hoja)
\begin{tikzpicture}[remember picture,overlay]
  \node[anchor=north east,inner sep=0pt]
    at ([xshift=-1.5cm,yshift=-1.2cm]current page.north east)
    {\includegraphics[width=9.75cm]{logos/logo-pba}};
\end{tikzpicture}

\centering
\vspace*{2.58cm}
% Títulos superiores del logo IMPA
{\fontsize{20.74}{25}\selectfont\bfseries Escuela de Educación Secundaria Técnica\par}
\vspace{0.2cm}
{\fontsize{20.74}{25}\selectfont\bfseries N° 7 - Quilmes\par}
\vspace{0.5cm}
% Logo IMPA centrado y grande
\includegraphics[width=0.775\linewidth]{logos/logo-impa}\par
\vspace{0.4cm}
% Título inferior del logo IMPA
{\fontsize{20.74}{25}\selectfont\bfseries Taller Regional Quilmes\par}
\vspace{2cm}
\fbox{\parbox{0.8\linewidth}{\centering\vspace{1em}
{\fontsize{24.88}{30}\selectfont\bfseries EMPRENDIMIENTOS\\ PRODUCTIVOS Y\\ DESARROLLO\\ LOCAL\par}\vspace{1em}}}\par
\vspace{0.8cm}
% Año / curso
{\fontsize{20.74}{25}\selectfont\bfseries 7mo año\par}
\vfill
{\fontsize{14.4}{17.3}\selectfont\bfseries EDICIÓN\par}
\vspace{0.2cm}
{\fontsize{14.4}{17.3}\selectfont\bfseries CICLO LECTIVO 2026\par}
\end{titlepage}
 % ============================================================
% Proyecto Pedagógico — páginas 2 a 5 del original, transcriptas
% a LaTeX tal cual (texto literal, mismo formato y cortes de página).
% Requiere en el preámbulo:
%   \usepackage{xcolor} \usepackage{fancyhdr} \usepackage{enumitem}
%   \usepackage{microtype} \usepackage{geometry}
% Va a continuación de la carátula: % ============================================================
% Carátula — reemplaza el bloque \begin{titlepage}...\end{titlepage}
% Requiere en el preámbulo: \usepackage{graphicx} y \usepackage{tikz}
% Logos en la carpeta logos/ (logo-impa.pdf y logo-pba.pdf)
% ============================================================
\begin{titlepage}
% Logo PBA en el margen superior derecho (posición absoluta en la hoja)
\begin{tikzpicture}[remember picture,overlay]
  \node[anchor=north east,inner sep=0pt]
    at ([xshift=-1.5cm,yshift=-1.2cm]current page.north east)
    {\includegraphics[width=9.75cm]{logos/logo-pba}};
\end{tikzpicture}

\centering
\vspace*{2.58cm}
% Títulos superiores del logo IMPA
{\fontsize{20.74}{25}\selectfont\bfseries Escuela de Educación Secundaria Técnica\par}
\vspace{0.2cm}
{\fontsize{20.74}{25}\selectfont\bfseries N° 7 - Quilmes\par}
\vspace{0.5cm}
% Logo IMPA centrado y grande
\includegraphics[width=0.775\linewidth]{logos/logo-impa}\par
\vspace{0.4cm}
% Título inferior del logo IMPA
{\fontsize{20.74}{25}\selectfont\bfseries Taller Regional Quilmes\par}
\vspace{2cm}
\fbox{\parbox{0.8\linewidth}{\centering\vspace{1em}
{\fontsize{24.88}{30}\selectfont\bfseries EMPRENDIMIENTOS\\ PRODUCTIVOS Y\\ DESARROLLO\\ LOCAL\par}\vspace{1em}}}\par
\vspace{0.8cm}
% Año / curso
{\fontsize{20.74}{25}\selectfont\bfseries 7mo año\par}
\vfill
{\fontsize{14.4}{17.3}\selectfont\bfseries EDICIÓN\par}
\vspace{0.2cm}
{\fontsize{14.4}{17.3}\selectfont\bfseries CICLO LECTIVO 2026\par}
\end{titlepage}
 % ============================================================
% Proyecto Pedagógico — páginas 2 a 5 del original, transcriptas
% a LaTeX tal cual (texto literal, mismo formato y cortes de página).
% Requiere en el preámbulo:
%   \usepackage{xcolor} \usepackage{fancyhdr} \usepackage{enumitem}
%   \usepackage{microtype} \usepackage{geometry}
% Va a continuación de la carátula: % ============================================================
% Carátula — reemplaza el bloque \begin{titlepage}...\end{titlepage}
% Requiere en el preámbulo: \usepackage{graphicx} y \usepackage{tikz}
% Logos en la carpeta logos/ (logo-impa.pdf y logo-pba.pdf)
% ============================================================
\begin{titlepage}
% Logo PBA en el margen superior derecho (posición absoluta en la hoja)
\begin{tikzpicture}[remember picture,overlay]
  \node[anchor=north east,inner sep=0pt]
    at ([xshift=-1.5cm,yshift=-1.2cm]current page.north east)
    {\includegraphics[width=9.75cm]{logos/logo-pba}};
\end{tikzpicture}

\centering
\vspace*{2.58cm}
% Títulos superiores del logo IMPA
{\fontsize{20.74}{25}\selectfont\bfseries Escuela de Educación Secundaria Técnica\par}
\vspace{0.2cm}
{\fontsize{20.74}{25}\selectfont\bfseries N° 7 - Quilmes\par}
\vspace{0.5cm}
% Logo IMPA centrado y grande
\includegraphics[width=0.775\linewidth]{logos/logo-impa}\par
\vspace{0.4cm}
% Título inferior del logo IMPA
{\fontsize{20.74}{25}\selectfont\bfseries Taller Regional Quilmes\par}
\vspace{2cm}
\fbox{\parbox{0.8\linewidth}{\centering\vspace{1em}
{\fontsize{24.88}{30}\selectfont\bfseries EMPRENDIMIENTOS\\ PRODUCTIVOS Y\\ DESARROLLO\\ LOCAL\par}\vspace{1em}}}\par
\vspace{0.8cm}
% Año / curso
{\fontsize{20.74}{25}\selectfont\bfseries 7mo año\par}
\vfill
{\fontsize{14.4}{17.3}\selectfont\bfseries EDICIÓN\par}
\vspace{0.2cm}
{\fontsize{14.4}{17.3}\selectfont\bfseries CICLO LECTIVO 2026\par}
\end{titlepage}
 % ============================================================
% Proyecto Pedagógico — páginas 2 a 5 del original, transcriptas
% a LaTeX tal cual (texto literal, mismo formato y cortes de página).
% Requiere en el preámbulo:
%   \usepackage{xcolor} \usepackage{fancyhdr} \usepackage{enumitem}
%   \usepackage{microtype} \usepackage{geometry}
% Va a continuación de la carátula: \input{caratula} \input{proyecto}
% ============================================================
\definecolor{ppAzul}{rgb}{0.267,0.447,0.769}
\definecolor{ppAzulClaro}{rgb}{0.557,0.667,0.859}

% Encabezado y pie del proyecto pedagógico
\fancypagestyle{proyecto}{%
  \fancyhf{}%
  \fancyhead[L]{\color{ppAzul}\bfseries\itshape 7mo Año}%
  \fancyhead[R]{\color{ppAzul}\bfseries\itshape Emprendimientos Productivos}%
  \fancyfoot[L]{\small\color{ppAzulClaro}\itshape scicolonejulian@impatrq.com}%
  \fancyfoot[C]{\raisebox{-0.7\height}{\includegraphics[height=0.8cm]{logos/logo-impa}}}% logo IMPA en el centro del pie
  \fancyfoot[R]{\small\color{ppAzulClaro}\itshape\thepage}%
  \renewcommand{\headrule}{{\color{ppAzul}\hrule height 0.5pt}}%
  \renewcommand{\footrule}{{\color{ppAzul}\hrule height 0.5pt}\vspace{2pt}}%
}

% Estilos del original
\newcommand{\ppTituloEsp}[1]{\underline{\textls[180]{\upshape\mdseries #1}}}      % OBJETIVOS..., espaciado y subrayado
\newcommand{\ppTitulo}[1]{\underline{\textbf{\textls[35]{#1}}}}                   % títulos en negrita cursiva subrayados
\newcommand{\ppNeg}[1]{\textbf{\textls[35]{#1}}}                                  % negrita cursiva espaciada
\newlist{guiones}{itemize}{1}
\setlist[guiones]{label=-,leftmargin=36pt,labelsep=18pt,itemsep=0pt,topsep=4pt,parsep=0pt}

\clearpage
\newgeometry{left=2.54cm,right=2.54cm,top=1.8cm,bottom=2cm,
             headheight=14pt,headsep=0.5cm,footskip=1.1cm}
\pagestyle{proyecto}
\setcounter{page}{2}
\begingroup
\fontfamily{lmr}\selectfont\itshape
\setlength{\parindent}{0pt}
\setlength{\parskip}{0pt}
\raggedright
\linespread{1.12}\selectfont

% ---------------------------------------------------------- página 2
\ppTituloEsp{FUNDAMENTACIÓN PEDAGOGICA:}

\vspace{8pt}
El Ciclo superior de la Secundaria técnica representa para los alumnos la oportunidad de
profundizar contenidos específicos de su orientación y prepararlos en un espacio de construcción
de nuevos conceptos.\\
Desde el presente espacio se pretende que los estudiantes comprendan las transformaciones
significativas en los últimos años, que van desde un nuevo contexto macroeconómico, vinculados
a la producción, la creación de empleo, el mejoramiento relativo de los indicadores sociales, la
generación de nuevas modalidades de gestión social y de emprendimientos productivos.\\
En ese sentido, el contexto actual, incentiva el debate en torno a una nueva conceptualización
del desarrollo, que se defina como inclusivo e integral. Es importante que sea inclusivo, ya que
implica una apuesta valorativa a la reconstrucción del tejido social y del mercado de trabajo,
tanto en relación a la creación de empleo como a la calidad del mismo. Es integral, ya que desde
un carácter multidimensional incorpora también lo social, lo político y lo cultural, revaloriza lo
ambiental y el territorio. Apuesta a una complementariedad sustentable entre el sector primario,
la industria y el sector servicios y busca asumir un carácter concertador entre el Estado, el
sector productivo, el social y el sector del conocimiento, para la definición del rumbo. Desde esta
perspectiva los emprendimientos productivos y el desarrollo local plantean una posibilidad y
alternativa importante para pensar el nuevo modelo. El desarrollo local pensado no solo en
relación al ámbito municipal, sino que cada vez más significa pensar el territorio, desde una
visión integral y flexible, que considere las diferentes escalas, que van desde lo local, micro
regional, lo interprovincial, e incluso macro regional, desde una perspectiva del MERCOSUR y
sus relaciones con el ámbito nacional y global y al mismo tiempo la trama institucional y de
actores existentes.

\vspace{22pt}
\ppTituloEsp{EXPECTATIVAS DE LOGRO:}

\vspace{10pt}
\hspace*{3pt}- Utilizar conceptos, herramientas, métodos y técnicas correspondientes a la gestión organizacional de
una empresa del ámbito aeronáutico en cualquiera de sus formas.

\vspace{8pt}
\hspace*{3pt}– Utilizar conceptos, herramientas, métodos y técnicas correspondientes a la gestión de los
recursos humanos

\vspace{8pt}
\hspace*{3pt}- Utilizar conceptos, herramientas, métodos y técnicas correspondientes a la gestión de calidad e
impacto ambiental.

\vspace{8pt}
- Interrelacionar la gestión global y la gestión de sus actividades específicas.

\vspace{8pt}
- Conocer el proceso productivo y la organización de una empresa, incluyendo comercialización,
selección y abastecimiento de sistemas, componentes y partes aeronáuticas y la generación y
participación de emprendimientos.

\vspace{8pt}
-Planificar, programar y organizar la producción, la prestación de servicios y el mantenimiento.

\vspace{8pt}
- Identificar proyectos productivos y de servicio, evaluar su factibilidad técnico-económica,
programar su implementación y su gestión.

% ---------------------------------------------------------- página 3
\newpage
\vspace*{12pt}
\ppTituloEsp{OBJETIVOS DE ENSEÑANZA:}

\begin{guiones}
  \item Desarrollar en el estudiante capacidad de análisis crítico y razonamiento.
  \item Promover el trabajo autónomo de los alumnos
  \item Estimular el establecimiento, comprobación y validación de normativas vigente.
  \item Promover el trabajo personal y grupal, valorando los aportes individuales y colectivos
        para la construcción del conocimiento de los riesgos.
  \item Estimular la mejora de la terminología y el vocabulario específico en las organizaciones.
  \item Fomentar la utilización de la normativa vigente como material de consulta y ampliación
        de lo trabajado en clase.
  \item Escuchar, registrar y retomar los aportes de los alumnos durante la clase.
\end{guiones}

\vspace{30pt}
{\fontsize{14}{17}\selectfont\underline{\textbf{Contenidos}}\par}

\vspace{8pt}
\ppTitulo{Unidad I: Introducción a los emprendimientos y desarrollos Productivos}

\begin{guiones}[leftmargin=18pt,labelsep=12pt,labelindent=-11pt]
  \item Teorías del Emprendedorismo, antecedentes históricos y el desarrollo del sector en la etapa
        actual. Tipos de Emprendedorismo.
  \item Nociones de derecho, finanzas, marketing y calidad en la gestión de emprendimientos.
  \item Análisis de Casos de Éxito de Desarrollos Productivos en el ámbito Aeronáutico de las
        siguientes empresas:
\end{guiones}

\vspace{8pt}
\hspace*{7pt}\ppNeg{FADEA (Fabrica Argentina de Aviones)}

\vspace{8pt}
\hspace*{8pt}\ppNeg{CICARE HELICOPTEROS S.A}

\vspace{8pt}
\hspace*{8pt}\ppNeg{REDIMEC}\\
\hspace*{4pt}\ppNeg{PAPIZ TURBINE}.

\vspace{8pt}
- Laboratorio de ideas y oportunidades. Planeamiento de negocios para emprendedores del sector
aeronáutico. Incubadoras: social; cultural y tecnológica.

\vspace{8pt}
- Desarrollo local y territorio: clúster, cadenas de valor, locales y regionales. Polos tecnológicos.

\vspace{8pt}
- El emprendedor y la promoción del desarrollo económico local: la planificación estratégica
participativa, las agencias de desarrollo, las incubadoras de empresas, los micros
emprendimientos y emprendimientos familiares.

\vspace{8pt}
- El Emprendedorismo como generador de valor en el sector aeronáutico y en una economía en
vías de desarrollo. Desarrollo rural, sustentabilidad del modelo y cuestiones ambientales.
Cooperación y asociativismo intermunicipal, micro regiones y desarrollo regional

% ---------------------------------------------------------- página 4
\newpage
\vspace*{12pt}
\hspace*{18pt}\ppTitulo{Unidad II: Gestión y procesos Operativos empresariales.}

\begin{guiones}
  \item Normas de Gestión de Calidad \ppNeg{ISO:9001} y Normas medioambientales \ppNeg{ISO 14.001}.
  \item Organigramas Empresariales.
  \item Introducción a procesos operativos.
  \item \ppNeg{Sistema Impositivo Argentino}.
  \item Metodologías \ppNeg{Keizen} – Metodología \ppNeg{LEAN} manufacture. Análisis FODA
  \item Manejos y Operaciones Bancarias:\ \ Conocimiento de distintos Productos bancarios.
        Financiación y \underline{CrowFounding}.
  \item Nociones básicas de Comex: Courier - Aduanas y régimen tributario aduanero
        Terminología \ppNeg{INCOTERM (FOB, CIF, CIP, DAP, DPU}).
  \item Procesos de Compras: Solicitud de Presupuestos – Ordenes de Compra – Facturación.
\end{guiones}

\vspace{22pt}
\hspace*{18pt}{\small\ppTitulo{Unidad III: Integración al ámbito laboral.}}

\begin{guiones}
  \item Introducción de Recursos Humanos; Procesos de Selección – Test Psicológicos.\\
        \ppNeg{Búsquedas de empleos - Taller de Armado de CV}
  \item Integración con equipos de trabajos interdepartamentales.
  \item Liquidación Salarial y Beneficios – Negociaciones Salariales.
\end{guiones}

\vspace{22pt}
\ppTitulo{EVALUACIÓN:}\\
El proceso de evaluación será continuo, con Trabajos Prácticos – Exámenes Teóricos y
Coloquios donde se relacionen todos los aspectos de cada unidad de aprendizaje.\\
El sistema de aprobación de la asignatura será a través de las siguientes pautas:

\vspace{8pt}
\ppNeg{Clases Teóricas}, \ppNeg{Teórico-Prácticas} y \ppNeg{Trabajos Prácticos}.

\vspace{16pt}
\ppTitulo{Criterios de evaluación:}

\begin{guiones}
  \item Interés y curiosidad los emprendimientos y desarrollos productivos
  \item Respeto y tolerancia hacia los demás.
  \item Asistencia regular a clase.
  \item Participación en clase.
  \item Asimilación y aplicación de los conceptos trabajados.
\end{guiones}

\vspace{8pt}
\ppTitulo{Instrumentos de evaluación:}

\begin{guiones}[itemsep=8pt]
  \item Entrega en tiempo y forma de trabajos prácticos.
  \item Observación del desempeño cotidiano
  \item[-] Evaluaciones escritas de cada unidad trabajada.\vspace{-8pt}
  \item Coloquios donde tengan que explicar y los criterios tomados
\end{guiones}

% ---------------------------------------------------------- página 5
\newpage
\ppTitulo{BIBLIOGRAFIA:}

\vspace{8pt}
- Normas ISO:9001:2015\ \ - Normas ISO 14.001

\vspace{8pt}
- El Empresario Argentino – Aldo Ferrer

\vspace{8pt}
- La economía argentina desde sus origines hasta el siglo XXI– Aldo Ferrer.

\vspace{8pt}
-\ \ País Riesgoso- Dobry, Hernan | Zuchovicki, Claudio

{\fontsize{12}{15.5}\selectfont
\vspace{8pt}
- La economía de tu vida – Tomas Bulat

\vspace{8pt}
- Economía descubierta – Tomas Bulat

\vspace{8pt}
- Estamos como somos: Por qué los argentinos no tenemos el país que queremos –\\
Tomas Bulat.

\vspace{8pt}
- Padre rico, padre pobre", por Robert Kiyosaki

\vspace{14pt}
- ¿Quién se ha llevado mi queso?", por Spencer Johnson

\vspace{8pt}
- El método Lean Start-Up" por Eric Ries.

\vspace{8pt}
- Biblioteca Electrónica de Resoluciones de AFIP\par}

\endgroup
\clearpage
\restoregeometry
\pagestyle{fancy}

% ============================================================
\definecolor{ppAzul}{rgb}{0.267,0.447,0.769}
\definecolor{ppAzulClaro}{rgb}{0.557,0.667,0.859}

% Encabezado y pie del proyecto pedagógico
\fancypagestyle{proyecto}{%
  \fancyhf{}%
  \fancyhead[L]{\color{ppAzul}\bfseries\itshape 7mo Año}%
  \fancyhead[R]{\color{ppAzul}\bfseries\itshape Emprendimientos Productivos}%
  \fancyfoot[L]{\small\color{ppAzulClaro}\itshape scicolonejulian@impatrq.com}%
  \fancyfoot[C]{\raisebox{-0.7\height}{\includegraphics[height=0.8cm]{logos/logo-impa}}}% logo IMPA en el centro del pie
  \fancyfoot[R]{\small\color{ppAzulClaro}\itshape\thepage}%
  \renewcommand{\headrule}{{\color{ppAzul}\hrule height 0.5pt}}%
  \renewcommand{\footrule}{{\color{ppAzul}\hrule height 0.5pt}\vspace{2pt}}%
}

% Estilos del original
\newcommand{\ppTituloEsp}[1]{\underline{\textls[180]{\upshape\mdseries #1}}}      % OBJETIVOS..., espaciado y subrayado
\newcommand{\ppTitulo}[1]{\underline{\textbf{\textls[35]{#1}}}}                   % títulos en negrita cursiva subrayados
\newcommand{\ppNeg}[1]{\textbf{\textls[35]{#1}}}                                  % negrita cursiva espaciada
\newlist{guiones}{itemize}{1}
\setlist[guiones]{label=-,leftmargin=36pt,labelsep=18pt,itemsep=0pt,topsep=4pt,parsep=0pt}

\clearpage
\newgeometry{left=2.54cm,right=2.54cm,top=1.8cm,bottom=2cm,
             headheight=14pt,headsep=0.5cm,footskip=1.1cm}
\pagestyle{proyecto}
\setcounter{page}{2}
\begingroup
\fontfamily{lmr}\selectfont\itshape
\setlength{\parindent}{0pt}
\setlength{\parskip}{0pt}
\raggedright
\linespread{1.12}\selectfont

% ---------------------------------------------------------- página 2
\ppTituloEsp{FUNDAMENTACIÓN PEDAGOGICA:}

\vspace{8pt}
El Ciclo superior de la Secundaria técnica representa para los alumnos la oportunidad de
profundizar contenidos específicos de su orientación y prepararlos en un espacio de construcción
de nuevos conceptos.\\
Desde el presente espacio se pretende que los estudiantes comprendan las transformaciones
significativas en los últimos años, que van desde un nuevo contexto macroeconómico, vinculados
a la producción, la creación de empleo, el mejoramiento relativo de los indicadores sociales, la
generación de nuevas modalidades de gestión social y de emprendimientos productivos.\\
En ese sentido, el contexto actual, incentiva el debate en torno a una nueva conceptualización
del desarrollo, que se defina como inclusivo e integral. Es importante que sea inclusivo, ya que
implica una apuesta valorativa a la reconstrucción del tejido social y del mercado de trabajo,
tanto en relación a la creación de empleo como a la calidad del mismo. Es integral, ya que desde
un carácter multidimensional incorpora también lo social, lo político y lo cultural, revaloriza lo
ambiental y el territorio. Apuesta a una complementariedad sustentable entre el sector primario,
la industria y el sector servicios y busca asumir un carácter concertador entre el Estado, el
sector productivo, el social y el sector del conocimiento, para la definición del rumbo. Desde esta
perspectiva los emprendimientos productivos y el desarrollo local plantean una posibilidad y
alternativa importante para pensar el nuevo modelo. El desarrollo local pensado no solo en
relación al ámbito municipal, sino que cada vez más significa pensar el territorio, desde una
visión integral y flexible, que considere las diferentes escalas, que van desde lo local, micro
regional, lo interprovincial, e incluso macro regional, desde una perspectiva del MERCOSUR y
sus relaciones con el ámbito nacional y global y al mismo tiempo la trama institucional y de
actores existentes.

\vspace{22pt}
\ppTituloEsp{EXPECTATIVAS DE LOGRO:}

\vspace{10pt}
\hspace*{3pt}- Utilizar conceptos, herramientas, métodos y técnicas correspondientes a la gestión organizacional de
una empresa del ámbito aeronáutico en cualquiera de sus formas.

\vspace{8pt}
\hspace*{3pt}– Utilizar conceptos, herramientas, métodos y técnicas correspondientes a la gestión de los
recursos humanos

\vspace{8pt}
\hspace*{3pt}- Utilizar conceptos, herramientas, métodos y técnicas correspondientes a la gestión de calidad e
impacto ambiental.

\vspace{8pt}
- Interrelacionar la gestión global y la gestión de sus actividades específicas.

\vspace{8pt}
- Conocer el proceso productivo y la organización de una empresa, incluyendo comercialización,
selección y abastecimiento de sistemas, componentes y partes aeronáuticas y la generación y
participación de emprendimientos.

\vspace{8pt}
-Planificar, programar y organizar la producción, la prestación de servicios y el mantenimiento.

\vspace{8pt}
- Identificar proyectos productivos y de servicio, evaluar su factibilidad técnico-económica,
programar su implementación y su gestión.

% ---------------------------------------------------------- página 3
\newpage
\vspace*{12pt}
\ppTituloEsp{OBJETIVOS DE ENSEÑANZA:}

\begin{guiones}
  \item Desarrollar en el estudiante capacidad de análisis crítico y razonamiento.
  \item Promover el trabajo autónomo de los alumnos
  \item Estimular el establecimiento, comprobación y validación de normativas vigente.
  \item Promover el trabajo personal y grupal, valorando los aportes individuales y colectivos
        para la construcción del conocimiento de los riesgos.
  \item Estimular la mejora de la terminología y el vocabulario específico en las organizaciones.
  \item Fomentar la utilización de la normativa vigente como material de consulta y ampliación
        de lo trabajado en clase.
  \item Escuchar, registrar y retomar los aportes de los alumnos durante la clase.
\end{guiones}

\vspace{30pt}
{\fontsize{14}{17}\selectfont\underline{\textbf{Contenidos}}\par}

\vspace{8pt}
\ppTitulo{Unidad I: Introducción a los emprendimientos y desarrollos Productivos}

\begin{guiones}[leftmargin=18pt,labelsep=12pt,labelindent=-11pt]
  \item Teorías del Emprendedorismo, antecedentes históricos y el desarrollo del sector en la etapa
        actual. Tipos de Emprendedorismo.
  \item Nociones de derecho, finanzas, marketing y calidad en la gestión de emprendimientos.
  \item Análisis de Casos de Éxito de Desarrollos Productivos en el ámbito Aeronáutico de las
        siguientes empresas:
\end{guiones}

\vspace{8pt}
\hspace*{7pt}\ppNeg{FADEA (Fabrica Argentina de Aviones)}

\vspace{8pt}
\hspace*{8pt}\ppNeg{CICARE HELICOPTEROS S.A}

\vspace{8pt}
\hspace*{8pt}\ppNeg{REDIMEC}\\
\hspace*{4pt}\ppNeg{PAPIZ TURBINE}.

\vspace{8pt}
- Laboratorio de ideas y oportunidades. Planeamiento de negocios para emprendedores del sector
aeronáutico. Incubadoras: social; cultural y tecnológica.

\vspace{8pt}
- Desarrollo local y territorio: clúster, cadenas de valor, locales y regionales. Polos tecnológicos.

\vspace{8pt}
- El emprendedor y la promoción del desarrollo económico local: la planificación estratégica
participativa, las agencias de desarrollo, las incubadoras de empresas, los micros
emprendimientos y emprendimientos familiares.

\vspace{8pt}
- El Emprendedorismo como generador de valor en el sector aeronáutico y en una economía en
vías de desarrollo. Desarrollo rural, sustentabilidad del modelo y cuestiones ambientales.
Cooperación y asociativismo intermunicipal, micro regiones y desarrollo regional

% ---------------------------------------------------------- página 4
\newpage
\vspace*{12pt}
\hspace*{18pt}\ppTitulo{Unidad II: Gestión y procesos Operativos empresariales.}

\begin{guiones}
  \item Normas de Gestión de Calidad \ppNeg{ISO:9001} y Normas medioambientales \ppNeg{ISO 14.001}.
  \item Organigramas Empresariales.
  \item Introducción a procesos operativos.
  \item \ppNeg{Sistema Impositivo Argentino}.
  \item Metodologías \ppNeg{Keizen} – Metodología \ppNeg{LEAN} manufacture. Análisis FODA
  \item Manejos y Operaciones Bancarias:\ \ Conocimiento de distintos Productos bancarios.
        Financiación y \underline{CrowFounding}.
  \item Nociones básicas de Comex: Courier - Aduanas y régimen tributario aduanero
        Terminología \ppNeg{INCOTERM (FOB, CIF, CIP, DAP, DPU}).
  \item Procesos de Compras: Solicitud de Presupuestos – Ordenes de Compra – Facturación.
\end{guiones}

\vspace{22pt}
\hspace*{18pt}{\small\ppTitulo{Unidad III: Integración al ámbito laboral.}}

\begin{guiones}
  \item Introducción de Recursos Humanos; Procesos de Selección – Test Psicológicos.\\
        \ppNeg{Búsquedas de empleos - Taller de Armado de CV}
  \item Integración con equipos de trabajos interdepartamentales.
  \item Liquidación Salarial y Beneficios – Negociaciones Salariales.
\end{guiones}

\vspace{22pt}
\ppTitulo{EVALUACIÓN:}\\
El proceso de evaluación será continuo, con Trabajos Prácticos – Exámenes Teóricos y
Coloquios donde se relacionen todos los aspectos de cada unidad de aprendizaje.\\
El sistema de aprobación de la asignatura será a través de las siguientes pautas:

\vspace{8pt}
\ppNeg{Clases Teóricas}, \ppNeg{Teórico-Prácticas} y \ppNeg{Trabajos Prácticos}.

\vspace{16pt}
\ppTitulo{Criterios de evaluación:}

\begin{guiones}
  \item Interés y curiosidad los emprendimientos y desarrollos productivos
  \item Respeto y tolerancia hacia los demás.
  \item Asistencia regular a clase.
  \item Participación en clase.
  \item Asimilación y aplicación de los conceptos trabajados.
\end{guiones}

\vspace{8pt}
\ppTitulo{Instrumentos de evaluación:}

\begin{guiones}[itemsep=8pt]
  \item Entrega en tiempo y forma de trabajos prácticos.
  \item Observación del desempeño cotidiano
  \item[-] Evaluaciones escritas de cada unidad trabajada.\vspace{-8pt}
  \item Coloquios donde tengan que explicar y los criterios tomados
\end{guiones}

% ---------------------------------------------------------- página 5
\newpage
\ppTitulo{BIBLIOGRAFIA:}

\vspace{8pt}
- Normas ISO:9001:2015\ \ - Normas ISO 14.001

\vspace{8pt}
- El Empresario Argentino – Aldo Ferrer

\vspace{8pt}
- La economía argentina desde sus origines hasta el siglo XXI– Aldo Ferrer.

\vspace{8pt}
-\ \ País Riesgoso- Dobry, Hernan | Zuchovicki, Claudio

{\fontsize{12}{15.5}\selectfont
\vspace{8pt}
- La economía de tu vida – Tomas Bulat

\vspace{8pt}
- Economía descubierta – Tomas Bulat

\vspace{8pt}
- Estamos como somos: Por qué los argentinos no tenemos el país que queremos –\\
Tomas Bulat.

\vspace{8pt}
- Padre rico, padre pobre", por Robert Kiyosaki

\vspace{14pt}
- ¿Quién se ha llevado mi queso?", por Spencer Johnson

\vspace{8pt}
- El método Lean Start-Up" por Eric Ries.

\vspace{8pt}
- Biblioteca Electrónica de Resoluciones de AFIP\par}

\endgroup
\clearpage
\restoregeometry
\pagestyle{fancy}

% ============================================================
\definecolor{ppAzul}{rgb}{0.267,0.447,0.769}
\definecolor{ppAzulClaro}{rgb}{0.557,0.667,0.859}

% Encabezado y pie del proyecto pedagógico
\fancypagestyle{proyecto}{%
  \fancyhf{}%
  \fancyhead[L]{\color{ppAzul}\bfseries\itshape 7mo Año}%
  \fancyhead[R]{\color{ppAzul}\bfseries\itshape Emprendimientos Productivos}%
  \fancyfoot[L]{\small\color{ppAzulClaro}\itshape scicolonejulian@impatrq.com}%
  \fancyfoot[C]{\raisebox{-0.7\height}{\includegraphics[height=0.8cm]{logos/logo-impa}}}% logo IMPA en el centro del pie
  \fancyfoot[R]{\small\color{ppAzulClaro}\itshape\thepage}%
  \renewcommand{\headrule}{{\color{ppAzul}\hrule height 0.5pt}}%
  \renewcommand{\footrule}{{\color{ppAzul}\hrule height 0.5pt}\vspace{2pt}}%
}

% Estilos del original
\newcommand{\ppTituloEsp}[1]{\underline{\textls[180]{\upshape\mdseries #1}}}      % OBJETIVOS..., espaciado y subrayado
\newcommand{\ppTitulo}[1]{\underline{\textbf{\textls[35]{#1}}}}                   % títulos en negrita cursiva subrayados
\newcommand{\ppNeg}[1]{\textbf{\textls[35]{#1}}}                                  % negrita cursiva espaciada
\newlist{guiones}{itemize}{1}
\setlist[guiones]{label=-,leftmargin=36pt,labelsep=18pt,itemsep=0pt,topsep=4pt,parsep=0pt}

\clearpage
\newgeometry{left=2.54cm,right=2.54cm,top=1.8cm,bottom=2cm,
             headheight=14pt,headsep=0.5cm,footskip=1.1cm}
\pagestyle{proyecto}
\setcounter{page}{2}
\begingroup
\fontfamily{lmr}\selectfont\itshape
\setlength{\parindent}{0pt}
\setlength{\parskip}{0pt}
\raggedright
\linespread{1.12}\selectfont

% ---------------------------------------------------------- página 2
\ppTituloEsp{FUNDAMENTACIÓN PEDAGOGICA:}

\vspace{8pt}
El Ciclo superior de la Secundaria técnica representa para los alumnos la oportunidad de
profundizar contenidos específicos de su orientación y prepararlos en un espacio de construcción
de nuevos conceptos.\\
Desde el presente espacio se pretende que los estudiantes comprendan las transformaciones
significativas en los últimos años, que van desde un nuevo contexto macroeconómico, vinculados
a la producción, la creación de empleo, el mejoramiento relativo de los indicadores sociales, la
generación de nuevas modalidades de gestión social y de emprendimientos productivos.\\
En ese sentido, el contexto actual, incentiva el debate en torno a una nueva conceptualización
del desarrollo, que se defina como inclusivo e integral. Es importante que sea inclusivo, ya que
implica una apuesta valorativa a la reconstrucción del tejido social y del mercado de trabajo,
tanto en relación a la creación de empleo como a la calidad del mismo. Es integral, ya que desde
un carácter multidimensional incorpora también lo social, lo político y lo cultural, revaloriza lo
ambiental y el territorio. Apuesta a una complementariedad sustentable entre el sector primario,
la industria y el sector servicios y busca asumir un carácter concertador entre el Estado, el
sector productivo, el social y el sector del conocimiento, para la definición del rumbo. Desde esta
perspectiva los emprendimientos productivos y el desarrollo local plantean una posibilidad y
alternativa importante para pensar el nuevo modelo. El desarrollo local pensado no solo en
relación al ámbito municipal, sino que cada vez más significa pensar el territorio, desde una
visión integral y flexible, que considere las diferentes escalas, que van desde lo local, micro
regional, lo interprovincial, e incluso macro regional, desde una perspectiva del MERCOSUR y
sus relaciones con el ámbito nacional y global y al mismo tiempo la trama institucional y de
actores existentes.

\vspace{22pt}
\ppTituloEsp{EXPECTATIVAS DE LOGRO:}

\vspace{10pt}
\hspace*{3pt}- Utilizar conceptos, herramientas, métodos y técnicas correspondientes a la gestión organizacional de
una empresa del ámbito aeronáutico en cualquiera de sus formas.

\vspace{8pt}
\hspace*{3pt}– Utilizar conceptos, herramientas, métodos y técnicas correspondientes a la gestión de los
recursos humanos

\vspace{8pt}
\hspace*{3pt}- Utilizar conceptos, herramientas, métodos y técnicas correspondientes a la gestión de calidad e
impacto ambiental.

\vspace{8pt}
- Interrelacionar la gestión global y la gestión de sus actividades específicas.

\vspace{8pt}
- Conocer el proceso productivo y la organización de una empresa, incluyendo comercialización,
selección y abastecimiento de sistemas, componentes y partes aeronáuticas y la generación y
participación de emprendimientos.

\vspace{8pt}
-Planificar, programar y organizar la producción, la prestación de servicios y el mantenimiento.

\vspace{8pt}
- Identificar proyectos productivos y de servicio, evaluar su factibilidad técnico-económica,
programar su implementación y su gestión.

% ---------------------------------------------------------- página 3
\newpage
\vspace*{12pt}
\ppTituloEsp{OBJETIVOS DE ENSEÑANZA:}

\begin{guiones}
  \item Desarrollar en el estudiante capacidad de análisis crítico y razonamiento.
  \item Promover el trabajo autónomo de los alumnos
  \item Estimular el establecimiento, comprobación y validación de normativas vigente.
  \item Promover el trabajo personal y grupal, valorando los aportes individuales y colectivos
        para la construcción del conocimiento de los riesgos.
  \item Estimular la mejora de la terminología y el vocabulario específico en las organizaciones.
  \item Fomentar la utilización de la normativa vigente como material de consulta y ampliación
        de lo trabajado en clase.
  \item Escuchar, registrar y retomar los aportes de los alumnos durante la clase.
\end{guiones}

\vspace{30pt}
{\fontsize{14}{17}\selectfont\underline{\textbf{Contenidos}}\par}

\vspace{8pt}
\ppTitulo{Unidad I: Introducción a los emprendimientos y desarrollos Productivos}

\begin{guiones}[leftmargin=18pt,labelsep=12pt,labelindent=-11pt]
  \item Teorías del Emprendedorismo, antecedentes históricos y el desarrollo del sector en la etapa
        actual. Tipos de Emprendedorismo.
  \item Nociones de derecho, finanzas, marketing y calidad en la gestión de emprendimientos.
  \item Análisis de Casos de Éxito de Desarrollos Productivos en el ámbito Aeronáutico de las
        siguientes empresas:
\end{guiones}

\vspace{8pt}
\hspace*{7pt}\ppNeg{FADEA (Fabrica Argentina de Aviones)}

\vspace{8pt}
\hspace*{8pt}\ppNeg{CICARE HELICOPTEROS S.A}

\vspace{8pt}
\hspace*{8pt}\ppNeg{REDIMEC}\\
\hspace*{4pt}\ppNeg{PAPIZ TURBINE}.

\vspace{8pt}
- Laboratorio de ideas y oportunidades. Planeamiento de negocios para emprendedores del sector
aeronáutico. Incubadoras: social; cultural y tecnológica.

\vspace{8pt}
- Desarrollo local y territorio: clúster, cadenas de valor, locales y regionales. Polos tecnológicos.

\vspace{8pt}
- El emprendedor y la promoción del desarrollo económico local: la planificación estratégica
participativa, las agencias de desarrollo, las incubadoras de empresas, los micros
emprendimientos y emprendimientos familiares.

\vspace{8pt}
- El Emprendedorismo como generador de valor en el sector aeronáutico y en una economía en
vías de desarrollo. Desarrollo rural, sustentabilidad del modelo y cuestiones ambientales.
Cooperación y asociativismo intermunicipal, micro regiones y desarrollo regional

% ---------------------------------------------------------- página 4
\newpage
\vspace*{12pt}
\hspace*{18pt}\ppTitulo{Unidad II: Gestión y procesos Operativos empresariales.}

\begin{guiones}
  \item Normas de Gestión de Calidad \ppNeg{ISO:9001} y Normas medioambientales \ppNeg{ISO 14.001}.
  \item Organigramas Empresariales.
  \item Introducción a procesos operativos.
  \item \ppNeg{Sistema Impositivo Argentino}.
  \item Metodologías \ppNeg{Keizen} – Metodología \ppNeg{LEAN} manufacture. Análisis FODA
  \item Manejos y Operaciones Bancarias:\ \ Conocimiento de distintos Productos bancarios.
        Financiación y \underline{CrowFounding}.
  \item Nociones básicas de Comex: Courier - Aduanas y régimen tributario aduanero
        Terminología \ppNeg{INCOTERM (FOB, CIF, CIP, DAP, DPU}).
  \item Procesos de Compras: Solicitud de Presupuestos – Ordenes de Compra – Facturación.
\end{guiones}

\vspace{22pt}
\hspace*{18pt}{\small\ppTitulo{Unidad III: Integración al ámbito laboral.}}

\begin{guiones}
  \item Introducción de Recursos Humanos; Procesos de Selección – Test Psicológicos.\\
        \ppNeg{Búsquedas de empleos - Taller de Armado de CV}
  \item Integración con equipos de trabajos interdepartamentales.
  \item Liquidación Salarial y Beneficios – Negociaciones Salariales.
\end{guiones}

\vspace{22pt}
\ppTitulo{EVALUACIÓN:}\\
El proceso de evaluación será continuo, con Trabajos Prácticos – Exámenes Teóricos y
Coloquios donde se relacionen todos los aspectos de cada unidad de aprendizaje.\\
El sistema de aprobación de la asignatura será a través de las siguientes pautas:

\vspace{8pt}
\ppNeg{Clases Teóricas}, \ppNeg{Teórico-Prácticas} y \ppNeg{Trabajos Prácticos}.

\vspace{16pt}
\ppTitulo{Criterios de evaluación:}

\begin{guiones}
  \item Interés y curiosidad los emprendimientos y desarrollos productivos
  \item Respeto y tolerancia hacia los demás.
  \item Asistencia regular a clase.
  \item Participación en clase.
  \item Asimilación y aplicación de los conceptos trabajados.
\end{guiones}

\vspace{8pt}
\ppTitulo{Instrumentos de evaluación:}

\begin{guiones}[itemsep=8pt]
  \item Entrega en tiempo y forma de trabajos prácticos.
  \item Observación del desempeño cotidiano
  \item[-] Evaluaciones escritas de cada unidad trabajada.\vspace{-8pt}
  \item Coloquios donde tengan que explicar y los criterios tomados
\end{guiones}

% ---------------------------------------------------------- página 5
\newpage
\ppTitulo{BIBLIOGRAFIA:}

\vspace{8pt}
- Normas ISO:9001:2015\ \ - Normas ISO 14.001

\vspace{8pt}
- El Empresario Argentino – Aldo Ferrer

\vspace{8pt}
- La economía argentina desde sus origines hasta el siglo XXI– Aldo Ferrer.

\vspace{8pt}
-\ \ País Riesgoso- Dobry, Hernan | Zuchovicki, Claudio

{\fontsize{12}{15.5}\selectfont
\vspace{8pt}
- La economía de tu vida – Tomas Bulat

\vspace{8pt}
- Economía descubierta – Tomas Bulat

\vspace{8pt}
- Estamos como somos: Por qué los argentinos no tenemos el país que queremos –\\
Tomas Bulat.

\vspace{8pt}
- Padre rico, padre pobre", por Robert Kiyosaki

\vspace{14pt}
- ¿Quién se ha llevado mi queso?", por Spencer Johnson

\vspace{8pt}
- El método Lean Start-Up" por Eric Ries.

\vspace{8pt}
- Biblioteca Electrónica de Resoluciones de AFIP\par}

\endgroup
\clearpage
\restoregeometry
\pagestyle{fancy}

% ============================================================
\definecolor{ppAzul}{rgb}{0.267,0.447,0.769}
\definecolor{ppAzulClaro}{rgb}{0.557,0.667,0.859}

% Encabezado y pie del proyecto pedagógico
\fancypagestyle{proyecto}{%
  \fancyhf{}%
  \fancyhead[L]{\color{ppAzul}\bfseries\itshape 7mo Año}%
  \fancyhead[R]{\color{ppAzul}\bfseries\itshape Emprendimientos Productivos}%
  \fancyfoot[L]{\small\color{ppAzulClaro}\itshape scicolonejulian@impatrq.com}%
  \fancyfoot[C]{\raisebox{-0.7\height}{\includegraphics[height=0.8cm]{logos/logo-impa}}}% logo IMPA en el centro del pie
  \fancyfoot[R]{\small\color{ppAzulClaro}\itshape\thepage}%
  \renewcommand{\headrule}{{\color{ppAzul}\hrule height 0.5pt}}%
  \renewcommand{\footrule}{{\color{ppAzul}\hrule height 0.5pt}\vspace{2pt}}%
}

% Estilos del original
\newcommand{\ppTituloEsp}[1]{\underline{\textls[180]{\upshape\mdseries #1}}}      % OBJETIVOS..., espaciado y subrayado
\newcommand{\ppTitulo}[1]{\underline{\textbf{\textls[35]{#1}}}}                   % títulos en negrita cursiva subrayados
\newcommand{\ppNeg}[1]{\textbf{\textls[35]{#1}}}                                  % negrita cursiva espaciada
\newlist{guiones}{itemize}{1}
\setlist[guiones]{label=-,leftmargin=36pt,labelsep=18pt,itemsep=0pt,topsep=4pt,parsep=0pt}

\clearpage
\newgeometry{left=2.54cm,right=2.54cm,top=1.8cm,bottom=2cm,
             headheight=14pt,headsep=0.5cm,footskip=1.1cm}
\pagestyle{proyecto}
\setcounter{page}{2}
\begingroup
\fontfamily{lmr}\selectfont\itshape
\setlength{\parindent}{0pt}
\setlength{\parskip}{0pt}
\raggedright
\linespread{1.12}\selectfont

% ---------------------------------------------------------- página 2
\ppTituloEsp{FUNDAMENTACIÓN PEDAGOGICA:}

\vspace{8pt}
El Ciclo superior de la Secundaria técnica representa para los alumnos la oportunidad de
profundizar contenidos específicos de su orientación y prepararlos en un espacio de construcción
de nuevos conceptos.\\
Desde el presente espacio se pretende que los estudiantes comprendan las transformaciones
significativas en los últimos años, que van desde un nuevo contexto macroeconómico, vinculados
a la producción, la creación de empleo, el mejoramiento relativo de los indicadores sociales, la
generación de nuevas modalidades de gestión social y de emprendimientos productivos.\\
En ese sentido, el contexto actual, incentiva el debate en torno a una nueva conceptualización
del desarrollo, que se defina como inclusivo e integral. Es importante que sea inclusivo, ya que
implica una apuesta valorativa a la reconstrucción del tejido social y del mercado de trabajo,
tanto en relación a la creación de empleo como a la calidad del mismo. Es integral, ya que desde
un carácter multidimensional incorpora también lo social, lo político y lo cultural, revaloriza lo
ambiental y el territorio. Apuesta a una complementariedad sustentable entre el sector primario,
la industria y el sector servicios y busca asumir un carácter concertador entre el Estado, el
sector productivo, el social y el sector del conocimiento, para la definición del rumbo. Desde esta
perspectiva los emprendimientos productivos y el desarrollo local plantean una posibilidad y
alternativa importante para pensar el nuevo modelo. El desarrollo local pensado no solo en
relación al ámbito municipal, sino que cada vez más significa pensar el territorio, desde una
visión integral y flexible, que considere las diferentes escalas, que van desde lo local, micro
regional, lo interprovincial, e incluso macro regional, desde una perspectiva del MERCOSUR y
sus relaciones con el ámbito nacional y global y al mismo tiempo la trama institucional y de
actores existentes.

\vspace{22pt}
\ppTituloEsp{EXPECTATIVAS DE LOGRO:}

\vspace{10pt}
\hspace*{3pt}- Utilizar conceptos, herramientas, métodos y técnicas correspondientes a la gestión organizacional de
una empresa del ámbito aeronáutico en cualquiera de sus formas.

\vspace{8pt}
\hspace*{3pt}– Utilizar conceptos, herramientas, métodos y técnicas correspondientes a la gestión de los
recursos humanos

\vspace{8pt}
\hspace*{3pt}- Utilizar conceptos, herramientas, métodos y técnicas correspondientes a la gestión de calidad e
impacto ambiental.

\vspace{8pt}
- Interrelacionar la gestión global y la gestión de sus actividades específicas.

\vspace{8pt}
- Conocer el proceso productivo y la organización de una empresa, incluyendo comercialización,
selección y abastecimiento de sistemas, componentes y partes aeronáuticas y la generación y
participación de emprendimientos.

\vspace{8pt}
-Planificar, programar y organizar la producción, la prestación de servicios y el mantenimiento.

\vspace{8pt}
- Identificar proyectos productivos y de servicio, evaluar su factibilidad técnico-económica,
programar su implementación y su gestión.

% ---------------------------------------------------------- página 3
\newpage
\vspace*{12pt}
\ppTituloEsp{OBJETIVOS DE ENSEÑANZA:}

\begin{guiones}
  \item Desarrollar en el estudiante capacidad de análisis crítico y razonamiento.
  \item Promover el trabajo autónomo de los alumnos
  \item Estimular el establecimiento, comprobación y validación de normativas vigente.
  \item Promover el trabajo personal y grupal, valorando los aportes individuales y colectivos
        para la construcción del conocimiento de los riesgos.
  \item Estimular la mejora de la terminología y el vocabulario específico en las organizaciones.
  \item Fomentar la utilización de la normativa vigente como material de consulta y ampliación
        de lo trabajado en clase.
  \item Escuchar, registrar y retomar los aportes de los alumnos durante la clase.
\end{guiones}

\vspace{30pt}
{\fontsize{14}{17}\selectfont\underline{\textbf{Contenidos}}\par}

\vspace{8pt}
\ppTitulo{Unidad I: Introducción a los emprendimientos y desarrollos Productivos}

\begin{guiones}[leftmargin=18pt,labelsep=12pt,labelindent=-11pt]
  \item Teorías del Emprendedorismo, antecedentes históricos y el desarrollo del sector en la etapa
        actual. Tipos de Emprendedorismo.
  \item Nociones de derecho, finanzas, marketing y calidad en la gestión de emprendimientos.
  \item Análisis de Casos de Éxito de Desarrollos Productivos en el ámbito Aeronáutico de las
        siguientes empresas:
\end{guiones}

\vspace{8pt}
\hspace*{7pt}\ppNeg{FADEA (Fabrica Argentina de Aviones)}

\vspace{8pt}
\hspace*{8pt}\ppNeg{CICARE HELICOPTEROS S.A}

\vspace{8pt}
\hspace*{8pt}\ppNeg{REDIMEC}\\
\hspace*{4pt}\ppNeg{PAPIZ TURBINE}.

\vspace{8pt}
- Laboratorio de ideas y oportunidades. Planeamiento de negocios para emprendedores del sector
aeronáutico. Incubadoras: social; cultural y tecnológica.

\vspace{8pt}
- Desarrollo local y territorio: clúster, cadenas de valor, locales y regionales. Polos tecnológicos.

\vspace{8pt}
- El emprendedor y la promoción del desarrollo económico local: la planificación estratégica
participativa, las agencias de desarrollo, las incubadoras de empresas, los micros
emprendimientos y emprendimientos familiares.

\vspace{8pt}
- El Emprendedorismo como generador de valor en el sector aeronáutico y en una economía en
vías de desarrollo. Desarrollo rural, sustentabilidad del modelo y cuestiones ambientales.
Cooperación y asociativismo intermunicipal, micro regiones y desarrollo regional

% ---------------------------------------------------------- página 4
\newpage
\vspace*{12pt}
\hspace*{18pt}\ppTitulo{Unidad II: Gestión y procesos Operativos empresariales.}

\begin{guiones}
  \item Normas de Gestión de Calidad \ppNeg{ISO:9001} y Normas medioambientales \ppNeg{ISO 14.001}.
  \item Organigramas Empresariales.
  \item Introducción a procesos operativos.
  \item \ppNeg{Sistema Impositivo Argentino}.
  \item Metodologías \ppNeg{Keizen} – Metodología \ppNeg{LEAN} manufacture. Análisis FODA
  \item Manejos y Operaciones Bancarias:\ \ Conocimiento de distintos Productos bancarios.
        Financiación y \underline{CrowFounding}.
  \item Nociones básicas de Comex: Courier - Aduanas y régimen tributario aduanero
        Terminología \ppNeg{INCOTERM (FOB, CIF, CIP, DAP, DPU}).
  \item Procesos de Compras: Solicitud de Presupuestos – Ordenes de Compra – Facturación.
\end{guiones}

\vspace{22pt}
\hspace*{18pt}{\small\ppTitulo{Unidad III: Integración al ámbito laboral.}}

\begin{guiones}
  \item Introducción de Recursos Humanos; Procesos de Selección – Test Psicológicos.\\
        \ppNeg{Búsquedas de empleos - Taller de Armado de CV}
  \item Integración con equipos de trabajos interdepartamentales.
  \item Liquidación Salarial y Beneficios – Negociaciones Salariales.
\end{guiones}

\vspace{22pt}
\ppTitulo{EVALUACIÓN:}\\
El proceso de evaluación será continuo, con Trabajos Prácticos – Exámenes Teóricos y
Coloquios donde se relacionen todos los aspectos de cada unidad de aprendizaje.\\
El sistema de aprobación de la asignatura será a través de las siguientes pautas:

\vspace{8pt}
\ppNeg{Clases Teóricas}, \ppNeg{Teórico-Prácticas} y \ppNeg{Trabajos Prácticos}.

\vspace{16pt}
\ppTitulo{Criterios de evaluación:}

\begin{guiones}
  \item Interés y curiosidad los emprendimientos y desarrollos productivos
  \item Respeto y tolerancia hacia los demás.
  \item Asistencia regular a clase.
  \item Participación en clase.
  \item Asimilación y aplicación de los conceptos trabajados.
\end{guiones}

\vspace{8pt}
\ppTitulo{Instrumentos de evaluación:}

\begin{guiones}[itemsep=8pt]
  \item Entrega en tiempo y forma de trabajos prácticos.
  \item Observación del desempeño cotidiano
  \item[-] Evaluaciones escritas de cada unidad trabajada.\vspace{-8pt}
  \item Coloquios donde tengan que explicar y los criterios tomados
\end{guiones}

% ---------------------------------------------------------- página 5
\newpage
\ppTitulo{BIBLIOGRAFIA:}

\vspace{8pt}
- Normas ISO:9001:2015\ \ - Normas ISO 14.001

\vspace{8pt}
- El Empresario Argentino – Aldo Ferrer

\vspace{8pt}
- La economía argentina desde sus origines hasta el siglo XXI– Aldo Ferrer.

\vspace{8pt}
-\ \ País Riesgoso- Dobry, Hernan | Zuchovicki, Claudio

{\fontsize{12}{15.5}\selectfont
\vspace{8pt}
- La economía de tu vida – Tomas Bulat

\vspace{8pt}
- Economía descubierta – Tomas Bulat

\vspace{8pt}
- Estamos como somos: Por qué los argentinos no tenemos el país que queremos –\\
Tomas Bulat.

\vspace{8pt}
- Padre rico, padre pobre", por Robert Kiyosaki

\vspace{14pt}
- ¿Quién se ha llevado mi queso?", por Spencer Johnson

\vspace{8pt}
- El método Lean Start-Up" por Eric Ries.

\vspace{8pt}
- Biblioteca Electrónica de Resoluciones de AFIP\par}

\endgroup
\clearpage
\restoregeometry
\pagestyle{fancy}
